\subsection{METRIZABLE UNIFORM SPACES}

DEFINITION 2. A metric on a set X is said to be compatible with a uniformity U on X if the uniformity defined by the metric coincides with U.

A uniformity on a set X is said to be metrizable if there is a metric on X compatible with this uniformity. A uniform space is said to be metrizable if its uniformity is metrizable.

Distinct metrics can be compatible with the same uniformity; they are then equivalent (§ 1, no. 2, Definition 2).

THEOREM 1. A uniformity is metrizable if and only if it is Hausdorff and the filter of entourages of the uniformity has a countable base.

The condition is necessary, for (with the notation of no. 2) the entourages \(V_{1/n}\) ( \(n \geqslant 1\) ) form a base of the filter of entourages of the uniformity of a metric space.

The condition is sufficient, for, if it is satisfied, the uniformity under consideration is defined by a single pseudometric, by Proposition 2 of § 1, no. 4; since the uniformity is Hausdorff, this pseudometric is a metric.

COROLLARY 1. A Hausdorff uniformity defined by a countable family of pseudometrics is metrizable.

For if \((f_n)\) is a sequence of pseudometrics defining such a structure, the filter of entourages is generated by the countable family of sets \(\overline{f}_n^1([0, 1/m])\) , where \(m\) and \(n\) each run through the set of integers \(>0\) .

COROLLARY 2. Every countable product of metrizable uniform spaces is metrizable.

For such a space is Hausdorff and its uniformity has a countable fundamental system of entourages (Chapter II, § 2, no. 6).

